% Section 7 -- drafted by Claude (2026-10-21), revised after two reviews.
\section{Limitations}\label{sec:limitations}

\begin{itemize}
  \item Every offline result holds out one of only four public benchmarks with enough pairs, and each formative score is one draw of 8 or 9 pairs. My standard errors treat pairs as independent, so they say little about a new benchmark; with the four benchmarks as the unit, V5's gain over V4 has a standard error near 0.003, not 0.0013.
  \item The ratings and guesses were tested on public benchmarks whose results Claude may have seen in training. In the pre-check, its three-example guesses matched Rasch ability at 0.957 on MathArena, whose per-model results are published (0.749 for the single-example guesses V5 uses). The offline gains may therefore exceed those on new benchmarks, and so may the hosted gains if the hidden benchmarks are public too; I cannot check this.
  \item V4 and V5 depend on a proprietary model's paid API (\$10.64 billed for the V5 run). Replies vary between calls and are not logged on the platform, so hosted runs cannot be reproduced exactly. Scoring sends subjects' recorded fields and some item texts to Anthropic.
  \item In the main offline tests every item had a rating, and each guess used one of the subject's own items as the example, a rated one for 229 of the 232 guesses (one per subject and benchmark). On Codabench, ratings and guesses that miss the 4-second wait are lost for that prediction, and a guess's example can be an image-only task or a bare instruction; the 270 of MathArena's 636 problems that share their text with others were never rated offline.
  \item Formative feedback, which comes from the hidden test pool, informed V3's $\lambda = 0$, taken against the offline evidence, and the switch from V4 to V5, taken under a rule committed before the upload. Two scenarios of the analysis behind V1's $\lambda = 0.5$ were also modeled on formative feedback, although the choice held without them.
  \item The pooling gain assumes that each pair's items are split independently, which the rules do not state. The rating trust in Eq.~\eqref{eq:trust} assumes a difficulty variance of $\tau^2 = 4$, while the fitted Rasch difficulties have a variance of about 5.6 within a benchmark (Appendix~\ref{app:model}).
\end{itemize}
